\documentclass[runningheads]{llncs}

\usepackage{polyglossia}
\setdefaultlanguage{english}
\setotherlanguage{greek}

\usepackage{fontspec}
\setmainfont{Times New Roman}
\usepackage{graphicx}
\usepackage{booktabs} % For professional looking tables

\begin{document}

\title{Ασφαλής επικοινωνία μεταξύ σενσώρων σε περιβάλλον Διαδύκτιου των πραγμάτων (IoT) με χρήση ελαφριάς κρυπτογραφίας μέσου του ASCON-128}

\author{Κωνσταντίνος Γκύλλης} 

\authorrunning{K.Γκύλλης}

\institute{Τμήμα Μηχανικών Η/Υ και πληροφορικής(CEID)\\
University of Patras, 26504, Greece \\
\email{up1100526@ac.upatras.gr}}

\maketitle

\begin{abstract}
Με την άνοδο του Διαδικτύου των Πραγμάτων (IoT), οι συσκευές έχουν να αντιμετωπίσουν μια μεγάλη πρόκληση που είναι η μη-ασφαλής επικοινωνία, 
αφού οι πληροφορίες μπορούν να υποκλαπούν από κάποιον κακόβουλο χρήστη (δηλαδή επιθέσεις τύπου Man-In-the-Middle), η οποία οφείλεται στην 
έλλειψη πόρων. Τα παραδοσιακά συστήματα κρυπτογραφίας απαιτούν τη χρήση μεγάλου ποσοστού της υπολογιστικής ισχύος για την κρυπτογράφηση 
και την αποκρυπτογράφηση. Έτσι, στην παρούσα εργασία θα υλοποιήσουμε μια σειρά από διεργασίες (pipeline) κάνοντας 
χρήση του αλγορίθμου ASCON-128, προσομοιώνοντας μια ασφαλή επικοινωνία μεταξύ δύο κόμβων (nodes) βασισμένων στην πλατφόρμα ESP32. 
Πιο αναλυτικά, θα υλοποιηθεί ένας ολοκληρωμένος κύκλος επικοινωνίας μεταξύ των δύο κόμβων: από τη συλλογή δεδομένων του αισθητήρα και τη δημιουργία δυαδικών πακέτων (binary packetization), 
έως την πιστοποιημένη κρυπτογράφηση (authenticated encryption) και την ασφαλή μετάδοση μέσω πρωτοκόλλου UDP. Θα δείξουμε ότι ο ASCON-128 παρέχει μια βιώσιμη ισορροπία μεταξύ 
ισχυρής ασφάλειας (εμπιστευτικότητα και ακεραιότητα) και των αυστηρών απαιτήσεων ενέργειας και απόδοσης του ενσωματωμένου υλικού (embedded hardware).

\keywords{Ασφάλεια στο Διαδίκτυο των Πραγμάτων (IoT) \and Ελαφριά Κρυπτογραφία \and ASCON-128 \and AEAD \and ESP32}
\end{abstract}

\section{Introduction}

Το οικοσύστημα του Διαδικτύου των Πραγμάτων (IoT) βασίζεται στην απρόσκοπτη ανταλλαγή πληροφοριών μέσα από κατανεμημένα δίκτυα επικοινωνίας. Ωστόσο, το υλικό (hardware) 
πάνω στο οποίο βασίζονται οι κόμβοι επικοινωνίας συχνά στερείται των απαραίτητων υπολογιστικών πόρων για την αποδοτική εκτέλεση καθιερωμένων πρωτοκόλλων ασφαλείας, όπως το TLS/SSL. 
Ως αποτέλεσμα, η ελαφριά κρυπτογραφία (lightweight cryptography) αναδεικνύεται ως ένας κρίσιμος τομέας με τεράστια εφαρμογή σε τέτοιου είδους περιορισμένα περιβάλλοντα.
Ο αλγόριθμος ASCON-128, ο οποίος αναδείχθηκε νικητής στον διαγωνισμό του οργανισμού NIST για την Ελαφριά Κρυπτογραφία χάρη στις εξαιρετικά χαμηλές απαιτήσεις του σε πόρους, προσφέρει Πιστοποιημένη Κρυπτογράφηση με Σχετιζόμενα Δεδομένα (Authenticated Encryption with Associated Data - AEAD). 

\subsection{ Στόχοι εργασίας}
Η παρούσα εργασία περιγράφει λεπτομερώς την υλοποίηση του ASCON-128 και παραθέτει τους εξής στόχους: 

\begin{itemize}
    \item \textbf{Ασφαλής σχεδίαση:}Τη σχεδίαση ενός ασφαλούς μηχανισμού επικοινωνίας μεταξύ συσκευών IoT.
    \item \textbf{Πιστοιποιημένη κρυπτογραφία:}Την εφαρμογή πιστοποιημένης κρυπτογράφησης (Authenticated Encryption with Associated Data - AEAD).
   %finish this \item \textbf{Fullfil the musts of security:}Την προστασία της εμπιστευτικότητας και της ακεραιότητας των δεδομένων κατά τη μετάδοση.
   % \item Τ\textbf{Συλλογή Δεδομένων:}ην αξιολόγηση της απόδοσης της κρυπτογράφησης σε συσκευές με περιορισμένους πόρους.
\end{itemize}

\subsection{Απαιτήσεις Συστήματος}

Για την επίτευξη των παραπάνω στόχων και τον ορθό σχεδιασμό της αρχιτεκτονικής, είναι απαραίτητος ο καθορισμός των βασικών προδιαγραφών λειτουργίας και ασφάλειας. Οι θεμελιώδεις απαιτήσεις του προτεινόμενου συστήματος ορίζονται ως εξής:

\begin{itemize}
    \item \textbf{Συλλογή Δεδομένων:} Το σύστημα πρέπει να συλλέγει αξιόπιστα δεδομένα από τους συνδεδεμένους αισθητήρες.
    \item \textbf{Ασύρματη Μετάδοση:} Τα δεδομένα πρέπει να μεταδίδονται ασύρματα και απρόσκοπτα μεταξύ των κόμβων IoT.
    \item \textbf{Εμπιστευτικότητα:} Τα μεταδιδόμενα δεδομένα πρέπει να προστατεύονται από κακόβουλες υποκλοπές (eavesdropping) κατά τη διάρκεια της μετάδοσης.
    \item \textbf{Ακεραιότητα:} Το σύστημα οφείλει να ανιχνεύει άμεσα οποιαδήποτε αλλοίωση ή τροποποίηση των δεδομένων.
    \item \textbf{Αποδοτικότητα Πόρων:} Η προτεινόμενη λύση πρέπει να είναι πλήρως λειτουργική σε συσκευές με αυστηρά περιορισμένους πόρους (embedded devices).
    \item \textbf{Κεντρική Παρακολούθηση:} Το σύστημα πρέπει να επιτρέπει την ασφαλή λήψη και παρακολούθηση των δεδομένων σε έναν τελικό κόμβο (sink node/dashboard).
    \item \textbf{Χαμηλή Καθυστέρηση:} Η καθυστέρηση (latency) που εισάγεται από τις κρυπτογραφικές διεργασίες πρέπει να διατηρείται σε ελάχιστα επίπεδα, ώστε να μην επηρεάζεται η λειτουργία σε πραγματικό χρόνο.
\end{itemize}

\section{System Architecture}
Η αρχιτεκτονική η οποία επιλέχθηκε για το μοντέλο της εργασία είναι σχεδιασμένη για μεταφορές με χαμηλή καθυστέρηση και εμπεριέχει 2 κόμβους της μορφής Παραγωγός και Καταναλωτής:
\begin{itemize}
    \item \textbf{Κόμβος A (Producer):} Υπεύθυνος για την εξαφάληση των δεδομέων I2C/SPI και την κρυπτογράφηση με τον ASCON-128.
    \item \textbf{Κόμβος B (Sink):} Υπεύθυνος για την επιβεβαίωση της ετικέτας των δεδομένων, αποκρυπτογράφηση (decryption) και την καταγραφή των δεδομένων σε κάποιο εξωτερικό διακομιστή (server).
\end{itemize}



\section{Πακετοποίηση Δεδομένων}

Με σκοπό την ελαχιστοποίηση του μεγέθους του μηνύματος που πρέπει να κρυπτογραφηθεί, καθώς και της ενέργειας που καταναλώνεται κατά τη μετάδοση, 
τα δεδομένα οργανώνονται σε μια συμπαγή δυαδική μορφή, όπως παρουσιάζεται στον Πίνακα~\ref{tab_packet}.

\begin{table}[!ht]
\caption{Δυαδική δομή ωφέλιμου φορτίου (payload) για την τηλεμετρία (telemetry) του αισθητήρα.}
\label{tab_packet}
\centering
\setlength{\tabcolsep}{10pt} % Αύξηση της απόστασης μεταξύ των στηλών
\begin{tabular}{l l c l}
\toprule
\textbf{Πεδίο} & \textbf{Τύπος} & \textbf{Μέγεθος (Bytes)} & \textbf{Περιγραφή} \\ 
\midrule
SensorID    & uint8\_t  & 1 & Αναγνωριστικό Υλικού \\
Temperature & float     & 4 & Αριθμός Κινητής Υποδιαστολής (IEEE 754) \\
Timestamp   & uint32\_t & 4 & Χρονοσφραγίδα Unix / Αύξων Αριθμός \\ 
\bottomrule
\end{tabular}
\end{table}

\section{ASCON-128 Implementation}
The implementation utilizes the ASCON-128 variant to provide 128-bit security. The AEAD process ensures that Node B can detect if a packet has been tampered with by an adversary. If the 128-bit authentication tag does not match the recalculated value at the receiver, the packet is discarded, effectively mitigating injection and bit-flipping attacks.

\section{Conclusion}
This implementation proves that ASCON-128 is highly suitable for the ESP32 platform, providing enterprise-grade security with negligible impact on transmission latency. Future work will assess power consumption profiles during high-frequency sampling.

\end{document}